%%%%%%%%%%%%%%%%%%%%%%%%%%%%%%%%%%%%%%%%%%%%%%%%%%%%%
%			HLAVIČKA								%
%%%%%%%%%%%%%%%%%%%%%%%%%%%%%%%%%%%%%%%%%%%%%%%%%%%%%
\documentclass[openany,12pt]{memoir}

\usepackage[utf8]{inputenc} 
\usepackage[czech]{babel}
\usepackage[T1]{fontenc}
\usepackage[top=1.5cm, bottom=2cm, left=2cm, right=2cm]{geometry}  % --> NASTAVENÍ OKRAJŮ
\usepackage{fancyhdr}
\usepackage{graphicx}
%\usepackage{lmodern}
\usepackage{tgschola}  % --- Silnější font
\usepackage{xwatermark}
\usepackage{xcolor}
\usepackage{changepage}
\usepackage{newunicodechar}



%%%%%% Package na zpěvník
\usepackage[full]{leadsheets}%http://mirrors.nic.cz/tex-archive/macros/latex/contrib/leadsheets/leadsheets_en.pdf   --> dokumentace	
\definesongtitletemplate{empty}{} 
\setchords{
format = \bfseries,   %tučné akordy
minor = {mi},% 
input-notation = {german},%
output-notation = {german}%
}
\definesongtitletemplate{empty}{} 

\newlength{\drop}
\newwatermark[allpages,color=red!50,angle=0,scale=2, xpos=0,ypos=0]{\includegraphics[width=5cm]{obr/pozadi2.jpg}} %--> dvojka na pozadí


%%%% Vlastní příkazy
\newcounter{Slokočet}   %Automatické číslování slok
\newcommand{\mezera}{\vspace*{0.5cm}}   %Horizontální odsazení slok
\newcommand{\stred}{5.2cm}   %%% Na zarovnání slok doprostřed, pozn. automatičtější zarovnávání na střed nejde
\newcommand{\refren}{\mezera \noindent \textbf{R:} } %refrén
\newcommand{\sloka}{\mezera \noindent \addtocounter{Slokočet}{1} \arabic{Slokočet}. } 	%sloka, která se automaticky čísluje
\newcommand{\ssloka}{\mezera \noindent}  % vlastní číslo sloky
\newcommand{\carka}{,\:}
\newcommand{\m}[1]{\color{white}{#1}}  %Pro akordy
\newcommand{\ap}{$'$} %Pro apostrof


%%%%%%%%%%%%%%%%%%%%%%%%%%%%%%%%%%%%%%%%%%%%%%%%%%%%%%
%			NÁVOD									 %
%%%%%%%%%%%%%%%%%%%%%%%%%%%%%%%%%%%%%%%%%%%%%%%%%%%%%%
%1. Věci v hlavičce IGNOROVAT
%2. Píseň psát do prostoru mezi \begin{song} a \end{song}
%3. další řádek se značí dvěma odsazeníma (= dvakrát stisknout enter)
%4. \refren vždy na začátku refrenu a \sloka na začátek sloky (automaticky se čísluje)
%  \ = alt gr + q ; [/] = alt gr f/g ; {/} = alt gr + b/n; ^ = alt gr + 3 + mezera
%Cokoliv napíšete do ^{  } se bude brát jako akord
%když se toto bude dotýkat nějakého slova (nebude mezi tím a slovem mezera)
%tak se akordy zjeví nad slovem, ale pište to před slovo
%Když se to nedotýká slova, tak akord lítá ve vzduchu a vytiskne se větší mezera
%První možnost je asi preferovanější
%5. Akordy stačí psát jen do první sloky, když se nezmění -- kytaristi to zvládnou



%\usepackage{subfiles}
%</++>

\begin{document}

\pagestyle{empty}
%\addtocontents{toc}{\protect\thispagestyle{empty}}
%\tableofcontents \thispagestyle{empty}\newpage
\newgeometry{top=1.5cm, bottom = 0cm, left = 2cm, right = 2cm}



%Pro kompilaci je důležitý, aby ten input a \end{document} zůstali na posledních
%dvou řádcích!!!!
\begin{song}{title=\centering Své Rocky Mountains \\\normalsize Vojta Kidák Tomáško  \vspace*{-0.3cm}}  %% sem se napíše jméno songu a autor
\moveright \stred \vbox{      %Varianta č. 1  ---> Jeden sloupec zarovnaný na střed	

\sloka 
	Je ^{A{\color{white}\_\_\_}}srpnový pátek a na horách sněží,

	^{D{\color{white}\_\_}}někdo jde vzhůru, někdo si věří,
	
	^{A{\color{white}\_\_}}krosna na zádech ^{D}tíží jak žulový ^{A{\color{white}\_}}blok.
	
	^{\phantom{.}}Zatím, co dole je léto a stromy
	
	a ^{D{\color{white}\_\_}}silácké řeči těch co se bojí,
	
	^{A}do bílé pláně ^{D{\color{white}\_\_}}udělat první ^{A{\color{white}\_}}krok.

\refren

\sloka
	^{D{\color{white}\_\_\_}}Vždycky se najdou a ^{G{\color{white}\_\_}}znovu to zkouší
	
	^{D{\color{white}\_\_}}slunce jak plamen ^{A{\color{white}\_}}bodá je v očích,
	
	^{E{\color{white}\_}}rvou se jak koně, ^{D}jen aby našli své ^{A}já.

\sloka	
	Sklonil jsem hlavu a ptal se sám sebe,
	
	jestli jak oni mám to svoje nebe
	
	tak blízko na dosah a nejsem na kolenou.
	
	Bez velkých řečí a ohraných frází,
	
	vstával jsem do tmy a nejmíň mi schází
	
	publikum, co chce mít za každou cenu své šou.

\refren

\sloka
	Možná mám poslední, poslední šanci,
	
	vzít svoji víru a v ledovém tanci
	
	na vrchol vynést svou vlajku a vidět ji vlát.
	
	Má každý před sebou své Rocky Mountains,
	
	svůj kopec z dětství, co zdál se být zlatý,
	
	ale kolik jich ztratilo cepín a muselo vzdát. 

\refren

\refren

\refren


}
\setcounter{Slokočet}{0}
\end{song}



\end{document}